\documentclass[11pt]{article}
\usepackage[utf8]{inputenc}
\usepackage[margin=1in]{geometry}
\usepackage{amsmath,amssymb,booktabs,bm}
\usepackage{graphicx}
\usepackage[colorlinks=true,linkcolor=blue]{hyperref}

\newcommand{\ii}{\mathrm{i}}
\newcommand{\rd}{\mathrm{d}}
\newcommand{\tr}{\mathrm{tr}}
\newcommand{\sgn}{\mathrm{sgn}}
\newcommand{\Id}{\mathbf{1}}
\newcommand{\TI}{\mathrm{TI}}
\newcommand{\triv}{\mathrm{triv}}
\newcommand{\strip}{\mathrm{strip}}

\begin{document}

\title{Slab-Geometry Note: Embedded TI/Trivial Interface and Local Chern Marker}
\author{}
\date{\today}

\maketitle

\section{Purpose}

This note records the setup and current numerical results for the slab/interface
calculation implemented in \texttt{src/interface\_chern}.  The geometry is a
finite explicit slab
\begin{equation}
  -M \le y \le M-1
\end{equation}
embedded between two semi-infinite bulk regions:
\begin{equation}
  y \ge M:\ \TI,\qquad y \le -M-1:\ \triv.
\end{equation}
Translation invariance is kept in the \(x\) direction, so the calculation is
done in mixed \((k_x,y)\) representation.

The main observable is the row-resolved local Chern-marker profile.  The two
results worth keeping as reference are:
\begin{itemize}
\item the equilibrium \(M=100\) slab develops a sharp interface crossover from a
  trivial-side marker near \(0\) to a TI-side plateau near \(-0.81\),
\item in the current monitored setup, a single monitored row at \(y=0\) already
  collapses the marker profile to nearly zero for all scanned \(\Gamma>0\).
\end{itemize}

\section{Geometry and Hamiltonian}

The bulk lattice Hamiltonian is the same Chern-insulator model used in
\texttt{monitored\_CI2.tex}:
\begin{equation}
  h_r(k_x,k_y)
  =
  t_0\Big[
    \sigma^1 \sin k_x
    + \sigma^2 \sin k_y
    + \sigma^3(r-\cos k_x-\cos k_y)
  \Big].
\end{equation}
We take
\begin{equation}
  r_{\TI}=1,\qquad r_{\triv}=3,
\end{equation}
so the upper bulk is topological and the lower bulk is trivial.

In mixed \((k_x,y)\) representation the explicit slab Hamiltonian is
\begin{equation}
  H_{\strip}(k_x)=
  \begin{pmatrix}
    h_{-M} & V & 0 & \cdots & 0\\
    V^\dagger & h_{-M+1} & V & \ddots & \vdots\\
    0 & V^\dagger & h_{-M+2} & \ddots & 0\\
    \vdots & \ddots & \ddots & \ddots & V\\
    0 & \cdots & 0 & V^\dagger & h_{M-1}
  \end{pmatrix},
\end{equation}
with
\begin{equation}
  h_y(k_x)
  =
  t_0\Big[
    \sigma^1\sin k_x
    + \sigma^3(r_y-\cos k_x)
  \Big],
\end{equation}
\begin{equation}
  V=\frac{\sigma^2}{2\ii}-\frac{\sigma^3}{2},
\end{equation}
and sharp interface profile
\begin{equation}
  r_y=
  \begin{cases}
    r_{\triv}, & -M \le y \le -1,\\[3pt]
    r_{\TI}, & 0 \le y \le M-1.
  \end{cases}
\end{equation}

\section{Lead Embedding}

The explicit slab is not given hard walls.  Instead, the boundary rows are
coupled to semi-infinite equilibrium leads through embedding self-energies:
\begin{equation}
  G_{\strip}^R(\omega,k_x)
  =
  \Big[
    (\omega+\ii \eta)\Id
    - H_{\strip}(k_x)
    - \Sigma_{\TI}^R(\omega,k_x)
    - \Sigma_{\triv}^R(\omega,k_x)
    - \Sigma_M^R
  \Big]^{-1}.
\end{equation}
The lead self-energies are obtained numerically from surface Green's functions
using Lopez-Sancho decimation.  In equilibrium, with no monitoring,
\begin{equation}
  G_{\strip}^K(\omega,k_x)
  =
  \sgn(\omega)\Big(G_{\strip}^R-G_{\strip}^A\Big),
\end{equation}
so the fixed bulks satisfy zero-temperature FDT exactly.

\section{Monitoring Setup}

The monitored runs use the following specific choice:
\begin{itemize}
\item only the single row \(y=0\) is monitored,
\item both \(\sigma^3\)-basis orbitals on that row are monitored,
\item the monitored retarded self-energy is
  \(\Sigma_M^R = -\ii \Gamma/2\) on the \(y=0\) row block,
\item for the large-\(M\) sweep we assume the half-filled monitored saddle,
  \(\Sigma_M^K=0\), and verify the corresponding monitored densities
  posteriori.
\end{itemize}

With this assumption the monitored calculation is a single pass rather than a
self-consistent iteration.  For the scanned \(M=100\) runs, the posterior
deviation of the monitored orbital densities from \(1/2\) is at most
\(4.7\times 10^{-11}\), so the half-filling assumption is numerically
self-consistent for this setup.

\section{Local Chern Marker}

The equal-time correlation matrix is
\begin{equation}
  C(k_x)
  =
  \frac{1}{2}\Id
  - \frac{\ii}{2}\int \frac{\rd\omega}{2\pi}\, G_{\strip}^K(\omega,k_x).
\end{equation}
The row-resolved marker is evaluated as
\begin{equation}
  c(y)
  =
  2\pi \int_{-\pi}^{\pi}\frac{\rd k_x}{2\pi}\,
  \tr_{\mathrm{orb}}
  \Big[
    \langle y|
    C(k_x)\,[\partial_{k_x}C(k_x),[Y,C(k_x)]]
    |y\rangle
  \Big],
\end{equation}
with
\begin{equation}
  Y_{yy'}=y\,\delta_{yy'}.
\end{equation}
For equilibrium runs this is the standard projector-based marker with
\(C(k_x)=P(k_x)\).  For monitored runs it is used as a nonequilibrium local
Hall/Chern-marker diagnostic.

\section{Alternative Boundary-First Reconstruction (\texorpdfstring{$M=0$}{M=0})}

An analytically cleaner route is to step back from the explicit slab and first
solve only the monitored \emph{boundary} problem at \(y=0\), then reconstruct
the Green's function deeper in the bulk from that boundary data.  This is most
natural for the semi-infinite edge geometry \(y\ge 0\), and is the right
starting point if the eventual goal is to build the equal-time correlation
matrix in a finite window near the edge.

Let \(g_{nm}(\omega,k_x)\), with \(n,m\ge 0\), denote the Green's function of
the \emph{unmonitored} half-space problem, and let the monitoring self-energy be
strictly local to the boundary row,
\begin{equation}
  \Sigma_M(\omega,k_x)
  =
  |0\rangle\,\hat\Sigma_M(\omega,k_x)\,\langle 0|,
\end{equation}
where \(\hat\Sigma_M\) is a \(2\times 2\) matrix in orbital space.  The full
Green's function satisfies the contour Dyson equation
\begin{equation}
  G = g + g\,\Sigma_M\,G.
\end{equation}
Because \(\Sigma_M\) has support only on \(y=0\), this equation closes exactly
into a boundary \(T\)-matrix problem:
\begin{equation}
  G_{nm} = g_{nm} + g_{n0}\,T\,g_{0m},
\end{equation}
with
\begin{equation}
  T
  =
  \bigl(1-\hat\Sigma_M g_{00}\bigr)^{-1}\hat\Sigma_M
  =
  \hat\Sigma_M\bigl(1-g_{00}\hat\Sigma_M\bigr)^{-1}.
\end{equation}

For the retarded and advanced components,
\begin{equation}
  G_{nm}^{R/A}
  =
  g_{nm}^{R/A}
  +
  g_{n0}^{R/A}\,T^{R/A}\,g_{0m}^{R/A},
\end{equation}
where
\begin{equation}
  T^{R/A}
  =
  \bigl(1-\hat\Sigma_M^{R/A} g_{00}^{R/A}\bigr)^{-1}\hat\Sigma_M^{R/A}.
\end{equation}
The Keldysh component follows from the Langreth rules,
\begin{equation}
  G_{nm}^{K}
  =
  g_{nm}^{K}
  +
  g_{n0}^{R}\,T^{R}\,g_{0m}^{K}
  +
  g_{n0}^{K}\,T^{A}\,g_{0m}^{A}
  +
  g_{n0}^{R}\,T^{K}\,g_{0m}^{A},
\end{equation}
with
\begin{equation}
  T^{K}
  =
  \bigl(1+T^{R}g_{00}^{R}\bigr)\hat\Sigma_M^{K}\bigl(1+g_{00}^{A}T^{A}\bigr)
  +
  T^{R}g_{00}^{K}T^{A}.
\end{equation}
Equivalently, all mixed and bulk blocks are reconstructed from the
unmonitored-half-space propagators \(g_{nm}^{R/A/K}\) and the local boundary
matrix \(T^{R/A/K}\).

Once \(G_{nm}^{K}\) is known, the equal-time correlation matrix in an edge
window \(0\le n,m\le N_{\mathrm{edge}}\) is
\begin{equation}
  C_{nm}(k_x)
  =
  \frac{1}{2}\delta_{nm}\Id_{\mathrm{orb}}
  -
  \frac{\ii}{2}\int\frac{\rd\omega}{2\pi}\,G_{nm}^{K}(\omega,k_x).
\end{equation}
This produces the full correlation matrix supported near the edge without ever
introducing a large explicit slab.

\paragraph{What information is needed beyond the current lead boundary condition?}

This point is important.  The present embedding code uses only the
\emph{surface-surface} Green's function \(g_{00}\) of the semi-infinite bulk in
order to build the self-energy.  That is sufficient to determine the local
boundary Green's function \(G_{00}\), but it is \emph{not} sufficient to
reconstruct \(G_{nm}\) in the bulk.  The formulas above show that one needs, in
addition,
\begin{equation}
  g_{n0}^{R/A/K},\qquad g_{0m}^{R/A/K},\qquad g_{nm}^{R/A/K}
  \quad (n,m>0),
\end{equation}
or, equivalently, the bulk propagation data that generates these objects.
Concretely, this means that the current lead boundary condition encodes only the
surface self-energy, while the boundary-first reconstruction requires
\emph{surface-to-bulk columns} and, for a whole edge window, the restricted
half-space block of the unmonitored bulk Green's function.

For a homogeneous nearest-neighbor bulk this extra information is still
calculable: it can be generated from the same half-space problem by a transfer
matrix or surface-to-bulk recursion.  The key warning is simply that the
surface self-energy alone does \emph{not} determine the bulk distribution of the
correlator.  One must supplement the current boundary condition with explicit
propagation data into the bulk.

\subsection*{Implementation-ready edge-window formulation}

For coding, the cleanest route is not to propagate directly from \(G_{00}\).
Instead, choose an edge window
\begin{equation}
  E=\{0,1,\dots,N_{\mathrm{edge}}\}
\end{equation}
inside the monitored half-space \(y\ge 0\), and compute the full Green's
function on that window exactly by embedding the \emph{tail} \(y\ge
N_{\mathrm{edge}}+1\) back into the last row of the window.

For the single-edge TI problem the half-space bulk is uniform, with
\begin{equation}
  h_{\mathrm{bulk}}(k_x)
  =
  t_0\Big[
    \sigma^1\sin k_x
    + \sigma^3(r_{\TI}-\cos k_x)
  \Big],
\end{equation}
\begin{equation}
  V=\frac{\sigma^2}{2\ii}-\frac{\sigma^3}{2}.
\end{equation}
The finite edge-window Hamiltonian is the block tridiagonal matrix
\begin{equation}
  H_E(k_x)=
  \begin{pmatrix}
    h_{\mathrm{bulk}} & V & 0 & \cdots & 0\\
    V^\dagger & h_{\mathrm{bulk}} & V & \ddots & \vdots\\
    0 & V^\dagger & h_{\mathrm{bulk}} & \ddots & 0\\
    \vdots & \ddots & \ddots & \ddots & V\\
    0 & \cdots & 0 & V^\dagger & h_{\mathrm{bulk}}
  \end{pmatrix},
\end{equation}
whose rows are indexed by \(n=0,\dots,N_{\mathrm{edge}}\).

Let \(g_{\mathrm{tail},s}^{R/A}(\omega,k_x)\) denote the surface Green's
function of the homogeneous semi-infinite tail \(n\ge N_{\mathrm{edge}}+1\).
It satisfies the usual surface Dyson equation
\begin{equation}
  g_{\mathrm{tail},s}^{R}
  =
  \Big[
    (\omega+\ii 0^+)\Id_{\mathrm{orb}}
    - h_{\mathrm{bulk}}
    - V g_{\mathrm{tail},s}^{R} V^\dagger
  \Big]^{-1},
\end{equation}
with \(g_{\mathrm{tail},s}^{A}=(g_{\mathrm{tail},s}^{R})^\dagger\).  The tail
self-energy acting on the last row of the edge window is
\begin{equation}
  \Sigma_{\mathrm{tail}}^{R/A/K}
  =
  |N_{\mathrm{edge}}\rangle\,
  \hat\Sigma_{\mathrm{tail}}^{R/A/K}\,
  \langle N_{\mathrm{edge}}|,
\end{equation}
\begin{equation}
  \hat\Sigma_{\mathrm{tail}}^{R/A}
  =
  V g_{\mathrm{tail},s}^{R/A} V^\dagger,
\end{equation}
\begin{equation}
  \hat\Sigma_{\mathrm{tail}}^{K}
  =
  \sgn(\omega)\Big(
    \hat\Sigma_{\mathrm{tail}}^{R}
    -
    \hat\Sigma_{\mathrm{tail}}^{A}
  \Big),
\end{equation}
since the bulk tail is fixed in equilibrium at zero temperature.

The exact \emph{unmonitored} Green's function on the entire edge window is then
\begin{equation}
  g_E^{R}
  =
  \Big[
    (\omega+\ii 0^+)\Id
    - H_E(k_x)
    - \Sigma_{\mathrm{tail}}^{R}
  \Big]^{-1},
\end{equation}
\begin{equation}
  g_E^{A}=(g_E^{R})^\dagger,
\end{equation}
\begin{equation}
  g_E^{K}
  =
  g_E^{R}\Sigma_{\mathrm{tail}}^{K}g_E^{A}
  =
  \sgn(\omega)\Big(g_E^{R}-g_E^{A}\Big).
\end{equation}
This gives all matrix elements \(g_{nm}^{R/A/K}\) for \(0\le n,m\le
N_{\mathrm{edge}}\) in one inversion, with no separate surface-to-bulk
recursion.

Monitoring is then added only on the boundary row \(n=0\):
\begin{equation}
  \Sigma_M^{R/A/K}
  =
  |0\rangle\,\hat\Sigma_M^{R/A/K}\,\langle 0|.
\end{equation}
For the monitored problem on the window one may either use the boundary
\(T\)-matrix formulas above, or simply solve the embedded window Dyson equation
directly,
\begin{equation}
  G_E^{R}
  =
  \Big[
    (\omega+\ii 0^+)\Id
    - H_E(k_x)
    - \Sigma_{\mathrm{tail}}^{R}
    - \Sigma_M^{R}
  \Big]^{-1},
\end{equation}
\begin{equation}
  G_E^{A}=(G_E^{R})^\dagger,
\end{equation}
\begin{equation}
  G_E^{K}
  =
  G_E^{R}\Big(
    \Sigma_{\mathrm{tail}}^{K}
    +
    \Sigma_M^{K}
  \Big)G_E^{A}.
\end{equation}
The correlation matrix supported near the edge is then
\begin{equation}
  C_{nm}(k_x)
  =
  \frac{1}{2}\delta_{nm}\Id_{\mathrm{orb}}
  - \frac{\ii}{2}\int\frac{\rd\omega}{2\pi}\,
  \bigl[G_E^{K}(\omega,k_x)\bigr]_{nm},
  \qquad 0\le n,m\le N_{\mathrm{edge}}.
\end{equation}

For the monitor used in the present code,
\begin{equation}
  \hat\Sigma_M^{R}=-\frac{\ii\Gamma}{2}\Id_{\mathrm{orb}},\qquad
  \hat\Sigma_M^{A}=+\frac{\ii\Gamma}{2}\Id_{\mathrm{orb}},
\end{equation}
and, if one keeps the general self-consistent form,
\begin{equation}
  \hat\Sigma_M^{K}
  =
  -\ii\Gamma
  \begin{pmatrix}
    1-2n_1 & 0\\
    0 & 1-2n_2
  \end{pmatrix},
\end{equation}
with monitored orbital occupancies extracted from the boundary block
\begin{equation}
  n_a
  =
  \int_{-\pi}^{\pi}\frac{\rd k_x}{2\pi}
  \left[
    \frac{1}{2}
    -
    \frac{\ii}{2}\int\frac{\rd\omega}{2\pi}\,
    \bigl[G_{00}^{K}(\omega,k_x)\bigr]_{aa}
  \right].
\end{equation}
At the half-filled monitored saddle used in the large-\(M\) scans,
\begin{equation}
  \hat\Sigma_M^{K}=0,
\end{equation}
so the full window Keldysh component reduces to
\begin{equation}
  G_E^{K}
  =
  G_E^{R}\Sigma_{\mathrm{tail}}^{K}G_E^{A}.
\end{equation}

\paragraph{Minimal data another agent actually needs.}

The implementation can therefore be based on the following minimal set of bulk
inputs:
\begin{enumerate}
\item the local bulk row Hamiltonian \(h_{\mathrm{bulk}}(k_x)\),
\item the nearest-row hopping \(V\),
\item the bulk-tail surface Green's function \(g_{\mathrm{tail},s}^{R/A}\),
\item equilibrium FDT for the tail, which fixes \(\Sigma_{\mathrm{tail}}^{K}\).
\end{enumerate}
No explicit bath distribution matrix beyond equilibrium FDT is needed.  What is
\emph{not} sufficient is the already-monitored local boundary Green's function
\(G_{00}\) by itself: one must still rebuild the whole edge window with the tail
embedding, because \(G_{00}\) alone does not determine the off-diagonal bulk
blocks \(G_{nm}\).

\section{Numerical Settings Used for the Reference Runs}

The large-\(M\) runs shown below use
\begin{equation}
  M=100,\qquad Nk_x=41,
\end{equation}
\begin{equation}
  \omega_{\mathrm{dense}}=2,\quad
  \omega_{\max}=6,\quad
  \Delta\omega_{\mathrm{dense}}=0.05,\quad
  \Delta\omega_{\mathrm{coarse}}=0.15.
\end{equation}
The code now enforces
\begin{equation}
  \eta \ge 2\Delta\omega_{\mathrm{dense}},
\end{equation}
so these runs use
\begin{equation}
  \eta=0.1.
\end{equation}
For the equilibrium \(M=100\) run the diagnostics are
\begin{equation}
  \texttt{fdt\_max\_error}=0,\qquad
  \texttt{lead\_residual\_max}=4.56\times 10^{-14}.
\end{equation}

\section{Results}

\subsection{Equilibrium slab}

The equilibrium \(M=100\) marker profile is shown in
Fig.~\ref{fig:eq-marker}.  Numerically,
\begin{equation}
  \bar c_{\TI}=-0.8066,\qquad \bar c_{\triv}=-0.00556,
\end{equation}
and representative rows are
\begin{equation}
  c(-1)=-0.1036,\qquad c(0)=-0.5511,\qquad c(1)=-0.8078.
\end{equation}
So the crossover is localized to the interface region, while the deep TI side
quickly saturates to a stable plateau.

\begin{figure}[t]
  \centering
  \includegraphics[width=0.82\linewidth]{output/interface_chern_M100_eq/interface_chern_M100_eq_marker.pdf}
  \caption{Equilibrium row-resolved Chern marker for the embedded
  \(M=100\) TI/trivial slab.  The trivial side stays near \(0\), while the TI
  side develops a plateau near \(-0.81\).}
  \label{fig:eq-marker}
\end{figure}

The corresponding equilibrium spectral output is shown in
Fig.~\ref{fig:eq-spectral}.  The top panel is the interface-window spectral sum
used by the code, while the remaining panels are selected row-resolved spectra
from the trivial side, interface region, and TI side.  This is the direct
spectral companion to the marker profile in Fig.~\ref{fig:eq-marker}.

\begin{figure}[p]
  \centering
  \includegraphics[width=\linewidth,height=0.78\textheight,keepaspectratio]{output/interface_chern_M100_eq/interface_chern_M100_eq_spectral.pdf}
  \caption{Equilibrium spectral output for the \(M=100\) slab.  The first panel
  is the interface-window sum; the remaining panels are row-resolved spectral
  functions for selected rows, with \(y<0\) in the left column and \(y\ge 0\) in
  the right column.  All row panels share one color scale.}
  \label{fig:eq-spectral}
\end{figure}

\subsection{Single-row monitoring at \texorpdfstring{$y=0$}{y=0}}

For a single monitored row \(y=0\), monitoring both \(\sigma^3\)-basis
orbitals, the marker profile is strongly suppressed.  At \(\Gamma=0.5\),
Fig.~\ref{fig:mon-marker} shows that the profile is already nearly flat, with
\begin{equation}
  \bar c_{\TI}=-1.21\times 10^{-4},\qquad
  \bar c_{\triv}=-2.16\times 10^{-5},
\end{equation}
and
\begin{equation}
  c(-1)\approx 0,\qquad c(0)\approx 0,\qquad c(1)\approx 0.
\end{equation}

\begin{figure}[t]
  \centering
  \includegraphics[width=0.82\linewidth]{output/interface_chern_M100_mon_half_fill/interface_chern_M100_mon_half_fill_marker.pdf}
  \caption{Monitored \(M=100\) marker profile for a single monitored row
  \(y=0\) with \(\Gamma=0.5\), monitoring both \(\sigma^3\)-basis orbitals and
  imposing the half-filled saddle \(\Sigma_M^K=0\).  The diagnostic collapses
  to nearly zero across the entire slab.}
  \label{fig:mon-marker}
\end{figure}

The monitored spectral output at the same parameters is shown in
Fig.~\ref{fig:mon-spectral}.  It is useful mainly as a reference diagnostic to
compare directly against the equilibrium spectral panels and against the marker
collapse in Fig.~\ref{fig:mon-marker}.

\begin{figure}[p]
  \centering
  \includegraphics[width=\linewidth,height=0.78\textheight,keepaspectratio]{output/interface_chern_M100_mon_half_fill/interface_chern_M100_mon_half_fill_spectral.pdf}
  \caption{Monitored spectral output for the \(M=100\) slab at \(\Gamma=0.5\),
  with only the row \(y=0\) monitored and the half-filled monitored saddle
  imposed.  As in the equilibrium case, the first panel is the interface-window
  sum, while the row-resolved spectral panels are arranged as \(y<0\) on the
  left and \(y\ge 0\) on the right with a shared row color scale.}
  \label{fig:mon-spectral}
\end{figure}

\subsection{Measurement-rate sweep}

The \(\Gamma\)-sweep is shown in Fig.~\ref{fig:gamma-sweep}.  The equilibrium
profile is included as the top-left reference panel, while the monitored panels
cover
\begin{equation}
  \Gamma = 0.05,\ 0.1,\ 0.2,\ 0.5,\ 1,\ 2,\ 5.
\end{equation}
Within this current setup, the collapse is essentially immediate: every
monitored panel is already close to zero for the full scanned range
\(\Gamma\ll t_0\) through \(\Gamma\gg t_0\).

This observation should be interpreted narrowly: it applies to the current
choice of monitoring both \(\sigma^3\) orbitals on a single row with the
half-filled monitored saddle enforced.  It is not yet evidence that any local
monitoring protocol would produce the same behavior.

\begin{figure}[p]
  \centering
  \includegraphics[width=\linewidth]{output/interface_chern_gamma_sweep_M100/interface_chern_gamma_sweep_M100_marker_panels_zoom.pdf}
  \caption{Panelized \(\Gamma\)-sweep for the \(M=100\) slab.  The equilibrium
  panel shows the nontrivial interface profile.  All monitored panels in the
  current setup are already collapsed near zero, even for \(\Gamma=0.05\ll
  t_0\).}
  \label{fig:gamma-sweep}
\end{figure}

\section{Takeaway}

For the embedded \(M=100\) slab, the equilibrium calculation shows a clear
interface-localized Chern-marker crossover.  In contrast, the present monitored
implementation --- single monitored row, both \(\sigma^3\) orbitals, half-filled
monitored saddle --- removes that profile almost completely for all scanned
nonzero \(\Gamma\).

The next natural follow-up calculations are:
\begin{itemize}
\item monitor only one orbital projector instead of both orbitals on the row,
\item repeat on smaller \(M\) with the full self-consistent \(\Sigma_M^K\),
\item compare the marker collapse against the corresponding spectral-function
  evolution at the interface.
\end{itemize}

\end{document}
